\begin{table}[!htbp]\centering
\caption{Main premium-enrollment outcomes: cubic-score OLS among high-math students}
\label{tab:high-math-premium-cubic-income-main}
\small
\begin{tabular}{lrrrr}\hline\hline
 & Institution & Field & Either & Both \\ \hline
Math: average derivative (per SD) & 15.60 (0.64) & 15.68 (0.56) & 12.86 (0.41) & 18.43 (0.65) \\
Verbal: average derivative (per SD) & 4.90 (0.32) & 1.51 (0.30) & 2.78 (0.19) & 3.64 (0.31) \\
Income decile 1 & Reference & Reference & Reference & Reference \\
Income decile 2 vs. 1 & 1.96 (1.36) & -0.23 (1.30) & 0.35 (1.16) & 1.39 (1.19) \\
Income decile 3 vs. 1 & 4.35 (1.28) & -2.62 (1.28) & -1.40 (1.14) & 3.12 (1.14) \\
Income decile 4 vs. 1 & 5.45 (1.28) & -2.30 (1.25) & -0.01 (1.15) & 3.16 (1.10) \\
Income decile 5 vs. 1 & 7.29 (1.27) & -2.40 (1.18) & 0.30 (1.06) & 4.58 (1.09) \\
Income decile 6 vs. 1 & 8.37 (1.27) & -2.28 (1.21) & 0.20 (1.04) & 5.89 (1.13) \\
Income decile 7 vs. 1 & 10.86 (1.29) & -3.29 (1.17) & 0.50 (1.06) & 7.08 (1.11) \\
Income decile 8 vs. 1 & 12.53 (1.31) & -3.72 (1.13) & 0.75 (0.99) & 8.06 (1.09) \\
Income decile 9 vs. 1 & 16.48 (1.42) & -5.13 (1.09) & 0.82 (0.95) & 10.54 (1.17) \\
Income decile 10 vs. 1 & 35.01 (1.86) & -17.24 (1.34) & 4.23 (0.98) & 13.54 (1.36) \\
\hline
Outcome mean (\%) & 52.63 & 62.97 & 80.78 & 34.82 \\
Full R-squared & 0.1667 & 0.0756 & 0.1308 & 0.0992 \\
R-squared without income & 0.1108 & 0.0573 & 0.1287 & 0.0927 \\
Income share of explained variation (\%) & 33.54 & 24.11 & 1.62 & 6.53 \\
Income indicators: joint clustered p-value & $<$0.001 & $<$0.001 & $<$0.001 & $<$0.001 \\
Observations & 93,423 & 93,423 & 93,423 & 93,423 \\
School clusters & 2,735 & 2,735 & 2,735 & 2,735 \\
Admission-year indicators & Yes & Yes & Yes & Yes \\
\hline\hline\end{tabular}
\begin{minipage}{0.98\textwidth}\footnotesize
\textit{Notes:} Income-decile coefficients and score average derivatives are in percentage points; school-clustered CR1 SEs are in parentheses. Score controls include $M,V,M^2,V^2,MV,M^3,V^3,M^2V,MV^2$. Income decile 1 is the reference. Admission-year effects are included in both full and reduced models. Income share is $(R^2_{full}-R^2_{without\ income})/R^2_{full}$, expressed as a percentage; it measures incremental explanatory power after scores and year, not a causal contribution. All models require original $z_{math} \geq 1$, without restandardizing scores or re-ranking income. The sample includes 8,522 non-enrollees and 840 admission-2026 records. Either is the inclusive union; all four outcomes are zero for non-enrollees.
\end{minipage}
\end{table}
